% الجزء: sets-functions-relations
% الفصل: relations
% القسم: special-properties

\documentclass[../../../include/open-logic-section]{subfiles}

\begin{document}

\olfileid[ar]{sfr}{rel}{prp}
\olsection{خصائص خاصة للعلاقات}

\begin{intro}
تبيّن أن بعض أنواع العلاقات شائعة إلى حد أنها مُنحت أسماء خاصة. فمثلًا،
تربط العلاقتان $\le$ و~$\subseteq$ عناصر مجاليهما (ولنقل $\Nat$ في حالة~$\le$،
و$\Pow{A}$ في حالة~$\subseteq$) بطرائق متشابهة. ولكي نحدد بدقة أوجه تشابه
هذه العلاقات وأوجه اختلافها، نصنفها بحسب بعض الخصائص الخاصة التي يمكن أن
تتصف بها العلاقات. ويتبيّن أن بعض هذه الخصائص (وبعض تراكيبها) مهم بوجه
خاص، ومنها الترتيبات وعلاقات التكافؤ.
\end{intro}

\begin{defn}[الانعكاسية]
تكون العلاقة $R \subseteq A^2$ \emph{انعكاسية} إذا وفقط إذا كان لكل $x \in
A$ أن $Rxx$ يصدق.
\end{defn}

\begin{defn}[التعدية]
تكون العلاقة $R \subseteq A^2$ \emph{متعدية} إذا وفقط إذا استلزم صدق $Rxy$
و$Ryz$ صدق $Rxz$ أيضًا.
\end{defn}

\begin{defn}[التناظر]
تكون العلاقة~$R \subseteq A^2$ \emph{متناظرة} إذا وفقط إذا استلزم صدق
$Rxy$ صدق~$Ryx$ أيضًا.
\end{defn}

\begin{defn}[مضادّة التناظر]
تكون العلاقة~$R \subseteq A^2$ \emph{مضادّة\- للتناظر\-} إذا وفقط إذا
استلزم صدق كل من $Rxy$ و$Ryx$ أن يكون $x=y$ (أو بعبارة أخرى: إذا كان $x\neq y$،
فلا بد أن يصدق واحد على الأقل من $\lnot Rxy$ و$\lnot Ryx$).
\end{defn}

\begin{explain}
في العلاقة المتناظرة، يصدق $Rxy$ و$Ryx$ معًا دائمًا، أو لا يصدق أي منهما.
وفي العلاقة المضادّة للتناظر، لا يمكن أن يصدق $Rxy$ و$Ryx$ معًا إلا إذا
كان $x = y$. لاحظ أن هذا لا \emph{يشترط} صدق $Rxy$ و$Ryx$ عندما يكون
$x = y$، بل يعني فقط أن صدقهما غير مستبعد. ولذلك يمكن أن تكون العلاقة
المضادّة للتناظر انعكاسية، ولكن ليست كل علاقة مضادّة للتناظر انعكاسية.
ولاحظ أيضًا أن كون العلاقة مضادّة للتناظر ومجرد كونها غير متناظرة شرطان
مختلفان. بل يمكن أن تكون العلاقة متناظرة ومضادّة للتناظر في الوقت نفسه
(ومثال ذلك علاقة الهوية).
\end{explain}

\begin{defn}[الاتصالية]
تكون العلاقة $R \subseteq A^2$ \emph{متصلة} إذا تحقق لكل $x,y\in
A$ أنه، متى كان $x \neq y$، يصدق واحد على الأقل من $Rxy$ و~$Ryx$.
\end{defn}

\begin{prob}
أعط أمثلة لعلاقات تكون (أ) انعكاسية ومتناظرة ولكن غير متعدية، و(ب)
انعكاسية ومضادّة للتناظر، و(ج) مضادّة للتناظر ومتعدية ولكن غير انعكاسية،
و(د) انعكاسية ومتناظرة ومتعدية. لا تستخدم علاقات على أعداد أو مجموعات.
\end{prob}

\begin{defn}[اللاانعكاسية]
تسمى العلاقة $R \subseteq A^2$ \emph{لا انعكاسية} إذا تحقق لكل $x \in
A$ أن $Rxx$ لا يصدق.
\end{defn}

\begin{defn}[اللاتناظرية]
تسمى العلاقة $R \subseteq A^2$ \emph{لا تناظرية} إذا لم يوجد زوج $x,y\in
A$ يصدق عنده كل من $Rxy$ و~$Ryx$.
\end{defn}

لاحظ أنه إذا كان $A \neq \emptyset$، فلا توجد علاقة لا انعكاسية على~$A$
تكون انعكاسية، وكل علاقة لا تناظرية على~$A$ مضادّة للتناظر أيضًا. ومع ذلك،
توجد علاقات $R \subseteq A^2$ لا تتصف بالانعكاسية ولا باللاانعكاسية، كما
توجد علاقات مضادّة للتناظر لا تتصف باللاتناظرية.

\end{document}
